\section{The program: prescribe the arithmetic evolution}\label{sec:intro}

The organizing proposal is to start with the exact shifted-zeta
differential equation and ask which field-dependent observable can
realize it. Loewner evolution supplies a controlled family of
contours; a supersymmetric defect theory supplies Wilson transport,
endpoint matter and candidate positive pairings. Their effective
boundary transfer is required to satisfy the arithmetic equation.
The shift equation thereby becomes a constraint on the choice of
observable and contour evolution, rather than a resemblance to be
tested only after those choices have been made.

This is an inverse realization problem. An equation for an operator
on boundary functions does not by itself determine a real Loewner
driver, a gauge connection or a quantum expectation. In particular,
the arithmetic generator is a nonlocal operator on an
infinite-dimensional function space, whereas a Wilson connection
initially acts on color indices along a curve. An explicit boundary
construction and an insertion identity must connect the two.

\subsection{The proposed chain of implications}

Write $V_{\omega,L}$ for the causal arithmetic transfer at shift
$\omega$ on an interval of length $L$. Its exact evolution and initial
condition are
\[
 \partial_\omega V_{\omega,L}=-G_{\omega,L}V_{\omega,L},
 \qquad V_{0,L}=I.
\]
The generator contains the gamma contribution, a fixed local term,
two pole terms and the prime-power translations. The objective is an
independently defined Wilson transfer $\mathcal V_{\omega,L}$
satisfying this equation, together with uniqueness in the relevant
causal class, so that $\mathcal V_{\omega,L}=V_{\omega,L}$.

A second identity should explain its norm:
\[
 \norm f^2=\norm{\mathcal V_{\omega,L}f}^2+
                 \norm{\mathcal B_{\omega,L}f}^2.
\]
Here $\mathcal B$ would be an independently constructed amplitude
map into a positive Hilbert space, potentially supplied by a
reflected Wilson network. This would prove contraction at the
parameters where the construction is valid. If those parameters
include $L_j\to\infty$ and $\omega_j\downarrow0$, the zero-shift
limit yields the full central Weil form on every compact support.
The precise conditional implication is
Proposition~\ref{prop:diagonal}.

There are therefore three substantive tasks: derive the prescribed
evolution from the observable, establish the cumulative norm
identity, and continue the construction to unbounded interval
length while the shift tends to zero. The known arithmetic
differential equation supplies a target for these tasks; it does
not discharge any of them.

\subsection{The role of the two papers}

This exposition retains the facts needed to assess that chain of
implications. Section~\ref{sec:arithmetic} states the full positivity
target, including the local constant $w_0$.
Section~\ref{sec:shift} derives the transfer and defect evolutions.
Sections~\ref{sec:realization}--\ref{sec:storage} formulate the
physical hypothesis and explain the constraints from the existing
calculations. Section~\ref{sec:fixed-window} gives the new fixed-window obstruction
and the localization controls, including their spatial limitations.
Section~\ref{sec:path} states the depth--shift continuation problem and
the separate fixed-window arithmetic progress, and Section~\ref{sec:outlook} gives the next
discriminating tasks.

The companion \emph{Supplementary information}~\cite{Supplement}
contains the detailed constructions, proofs and diagnostics.
A result that materially restricts the proposed realization is
summarized here even when its derivation is deferred there.
This division lets the main argument be read without the Clifford
algebra or perturbative integrals, while keeping those calculations
available as explicit tests of a candidate model.

\subsection{Sources and relation to earlier work}

The scalar-coupled Wilson framework originates in
Maldacena's construction~\cite{Maldacena}; tangent-dependent scalar
couplings and common line supersymmetry are due to
Zarembo~\cite{Zarembo}. The D3--D5 defect theory and its projector
conventions come from DeWolfe--Freedman--Ooguri and
Gaiotto--Witten~\cite{DFO,GW}. Baker's open Wilson lines and
semicircle are the endpoint controls~\cite{Baker}.
We use chordal Loewner normalization as in Lawler~\cite{Lawler}.
The Loewner/CFT relation of Bauer--Bernard~\cite{BauerBernard}
motivates looking for a closed averaged observable, but supplies
no Wilson realization of the arithmetic generator here.

Ordinary reflection and its distinction from an internal rotation
are guided by Osterwalder--Schrader and
Dorn--Pershin~\cite{OS,DP}. The compact-support Weil criterion is
classical; our convention agrees with Suzuki~\cite{Suzuki}.
The exact shift evolution is reproduced from the shared
background~\cite{SharedBackground}. The Wilson-lines, Loewner and
fractional-dimension investigations supply the previous realization
tests~\cite{ParentWilson,ParentLoewner,ParentFractional}.
The cumulative-storage continuation comes from the shifted-zeta
program~\cite{CumulativePath,AdaptivePath}. The new organization
follows Baker's proposal recorded in the shift-flow research
note~\cite{ShiftFlowNote}; it is a change in the research question,
not a claim to have solved the inverse problem.
